% ============================================================
%  Myora Academic Poster (A0 portrait)
%  Compile with: pdflatex -interaction=nonstopmode myora_academic_poster.tex
% ============================================================
\documentclass[final]{beamer}

\usepackage[orientation=portrait,size=a0,scale=1.45]{beamerposter}
\usepackage[utf8]{inputenc}
\usepackage[T1]{fontenc}
\usepackage[turkish]{babel}
\AtBeginDocument{\shorthandoff{=}}
\usepackage{lmodern}
\usepackage{microtype}
\usepackage{booktabs}
\usepackage{tabularx}
\usepackage{ragged2e}
\usepackage{tikz}
\usepackage{xcolor}
\usepackage{hyperref}
\usepackage{amsmath}
\usepackage{amssymb}

\usetikzlibrary{arrows.meta,positioning,shapes.geometric,shapes.misc,fit,calc,babel,backgrounds,chains}

\definecolor{myoraBlue}{HTML}{2563EB}
\definecolor{myoraDeep}{HTML}{0B2A63}
\definecolor{myoraCyan}{HTML}{06B6D4}
\definecolor{myoraGreen}{HTML}{10B981}
\definecolor{myoraPurple}{HTML}{7C3AED}
\definecolor{myoraOrange}{HTML}{F97316}
\definecolor{myoraRed}{HTML}{EF4444}
\definecolor{myoraYellow}{HTML}{F59E0B}
\definecolor{myoraInk}{HTML}{0F172A}
\definecolor{myoraMuted}{HTML}{475569}
\definecolor{myoraPanel}{HTML}{F8FAFC}
\definecolor{myoraLine}{HTML}{CBD5E1}
\definecolor{myoraSoft}{HTML}{E2E8F0}

\hypersetup{colorlinks=true,urlcolor=myoraBlue,linkcolor=myoraBlue}

\setbeamertemplate{navigation symbols}{}
\setbeamercolor{background canvas}{bg=white}
\setbeamercolor{block title}{fg=white,bg=myoraBlue}
\setbeamercolor{block body}{fg=myoraInk,bg=myoraPanel}
\setbeamerfont{block title}{size=\Large,series=\bfseries}
\setbeamerfont{block body}{size=\normalsize}
\setlength{\parindent}{0pt}
\setlength{\parskip}{0.45em}

\newcolumntype{Y}{>{\RaggedRight\arraybackslash}X}

% --------------------------------------------------------
% Helper macros
% --------------------------------------------------------

% Small coloured pill tag
\newcommand{\tagpill}[2]{%
  \tikz[baseline=(x.base)] \node[
    rounded corners=7pt,
    fill=#1!14,
    draw=#1!55,
    inner xsep=8pt,
    inner ysep=4pt,
    text=myoraInk
  ] (x) {\small #2};%
}

% Colored feature card (icon letter + title + description)
\newcommand{\featcard}[4]{%
  \begin{tikzpicture}
    \node[
      rounded corners=10pt,
      draw=#1!55, fill=#1!10,
      line width=1pt,
      minimum width=0.48\linewidth,
      minimum height=2.8cm,
      text width=0.44\linewidth,
      align=left,
      inner sep=8pt
    ] {%
      \textcolor{#1}{\Large\bfseries #2}\ \textbf{#3}\\[-0.15em]
      {\small #4}%
    };
  \end{tikzpicture}%
}

% Pipeline step box
\newcommand{\pipestep}[3]{%
  \node[
    rounded corners=8pt,
    draw=#1!60, fill=#1!12,
    line width=1pt,
    minimum height=2.05cm,
    text width=3.5cm,
    align=center,
    inner sep=5pt,
    font=\small
  ] (#2) {\textcolor{#1}{\small\bfseries #3}};%
}

% Data entity box
\newcommand{\entity}[2]{%
  \node[
    rounded corners=6pt,
    draw=myoraBlue!60,
    fill=myoraBlue!8,
    line width=0.9pt,
    minimum width=6.4cm,
    minimum height=1.4cm,
    align=center,
    font=\small\ttfamily
  ] (#1) {#2};%
}

% Horizontal bar (used for RDA chart)
\newcommand{\rdabar}[3]{%
  % #1=label, #2=percent (0-100), #3=color
  \begin{tikzpicture}[yscale=0.7]
    \draw[rounded corners=3pt, fill=myoraSoft, draw=myoraSoft]
      (0,0) rectangle (10,0.7);
    \draw[rounded corners=3pt, fill=#3, draw=#3]
      (0,0) rectangle ({#2/10},0.7);
    \node[anchor=east, font=\small\bfseries, text=myoraInk]
      at (-0.15,0.35) {#1};
    \node[anchor=west, font=\small, text=myoraMuted]
      at (10.2,0.35) {#2\%};
    \draw[dashed, thick, myoraGreen]
      (10,-0.05) -- (10,0.75);
  \end{tikzpicture}%
}

\newenvironment{posterblock}[1]{%
  \begin{block}{#1}
  \justifying
}{%
  \end{block}
}

\title{Myora: Yapay Zekâ Destekli Kişisel Sağlık ve Beslenme Takip Uygulaması}
\author{Seyit Ali Yorgun}
\institute{OSTİM Teknik Üniversitesi -- Yazılım Mühendisliği}
\date{2026}

\begin{document}
\begin{frame}[t]

% ============================================================
% HEADER BANNER
% ============================================================
\begin{tikzpicture}[remember picture,overlay]
  \fill[myoraDeep] (current page.north west) rectangle
    ($(current page.north east)+(0,-12cm)$);
  \fill[myoraCyan!80] ($(current page.north west)+(0,-12cm)$)
    rectangle ($(current page.north east)+(0,-12.4cm)$);

  % Logo circle on left
  \node[
    circle, draw=white, line width=2pt,
    minimum size=3.6cm, fill=myoraDeep,
    text=white
  ] (logo) at ($(current page.north west)+(4.0cm,-4.3cm)$)
    {\fontsize{38}{42}\selectfont\bfseries M};

  % Title block
  \node[anchor=north west, text=white, align=left]
    at ($(current page.north west)+(6.8cm,-1.6cm)$) {%
      {\fontsize{58}{62}\selectfont\bfseries MYORA}\\[0.2cm]
      {\fontsize{24}{28}\selectfont Yapay Zekâ Destekli Kişisel Sağlık ve Beslenme Takip Uygulaması}\\[0.45cm]
      \tikz\node[rounded corners=6pt,fill=myoraGreen!85,inner xsep=9pt,inner ysep=4pt,text=white]{\small\bfseries Bilgisayarlı Görü};\ %
      \tikz\node[rounded corners=6pt,fill=myoraOrange!85,inner xsep=9pt,inner ysep=4pt,text=white]{\small\bfseries Doğal Dil İşleme};\ %
      \tikz\node[rounded corners=6pt,fill=myoraPurple!85,inner xsep=9pt,inner ysep=4pt,text=white]{\small\bfseries Kişiselleştirilmiş RDA Motoru};
    };

  % Center subtitle area (lower half of banner)
  \node[anchor=center, text=white!95, align=center]
    at ($(current page.north)+(0,-9.0cm)$) {%
      {\large\bfseries Çok Platformlu Mobil Uygulama}\\[0.35cm]
      {\small React 18 + TypeScript\ \ \textbullet\ \ Capacitor 7\ \ \textbullet\ \ Supabase + Edge Functions\ \ \textbullet\ \ OpenAI GPT-4o (çok modlu)}
    };

  % Tech stack on the right
  \node[anchor=north east, text=white, align=right]
    at ($(current page.north east)+(-3.0cm,-1.2cm)$) {%
      {\normalsize\bfseries Teknoloji Yığını}\\[0.2cm]
      {\small \textbullet\ Android ve iOS (Capacitor 7)}\\[0.05cm]
      {\small \textbullet\ React 18 + TypeScript + Vite}\\[0.05cm]
      {\small \textbullet\ Supabase (PostgreSQL + Edge)}\\[0.05cm]
      {\small \textbullet\ OpenAI GPT-4o Vision API}\\[0.05cm]
      {\small \textbullet\ RAG (pgvector) + Zod}\\[0.05cm]
      {\small \textbullet\ TanStack Query + shadcn/ui}%
    };
\end{tikzpicture}

\vspace*{12.6cm}

% ============================================================
% COLUMNS
% ============================================================
\begin{columns}[t,totalwidth=\textwidth]

% ============================================================
% COLUMN 1
% ============================================================
\begin{column}{0.32\textwidth}

\begin{posterblock}{Özet}
\small
\textbf{Myora}, yapay zekâ bilgisayarlı görüsü, doğal dil işleme ve kişiselleştirilmiş bir besin motorunun birleşimiyle manuel sağlık kaydının önündeki engelleri kaldıran çok platformlu bir sağlık ve beslenme takip uygulamasıdır.

Kullanıcılar yemeklerini fotoğraflayarak anlık makro besin dökümü alabilir, sesli girdiyle günlük tutabilir, kan tahlili ve tıbbi belgeleri yapay zekâya analiz ettirebilir ve sonuçlara göre kişiselleştirilmiş takviye önerileri elde edebilir.

Sistem; Android ve iOS için \textbf{Capacitor} ile paketlenmiş bir React istemci, \textbf{Supabase} bulut altyapısı (PostgreSQL + Edge Functions) ve \textbf{OpenAI GPT-4o} çok modlu API üzerine kurulmuştur.
\end{posterblock}

\begin{posterblock}{Problem Tanımı}
\small
Geleneksel sağlık takip uygulamaları birçok sorundan muzdariptir:
\begin{itemize}\setlength{\itemsep}{0.15em}
  \item \textcolor{myoraRed}{\textbf{Yüksek manuel yük}}\,--\,kullanıcı her yiyeceği arayıp elle girmek zorundadır.
  \item \textcolor{myoraRed}{\textbf{Genel öneriler}}\,--\,herkese uyan tek kalori hedefi yaş, cinsiyet, VKİ, aktivite ve kan biyobelirteçlerini göz ardı eder.
  \item \textcolor{myoraRed}{\textbf{Parçalı veriler}}\,--\,beslenme, aktivite, belge ve giyilebilir cihazlar ayrı silolarda yaşar.
  \item \textcolor{myoraRed}{\textbf{Düşük sürdürülebilirlik}}\,--\,karmaşık arayüz, 2 hafta içinde uygulamayı bırakmaya yol açar.
\end{itemize}
\textbf{Myora'nın çözümü:} Manuel girişin yerini \emph{yakalama} alır\,--\,fotoğraf, ses veya belge yüklemesi\,--\,ve tüm sağlık sinyalleri tek bir kişiselleştirilmiş yapay zekâ katmanında birleşir.
\end{posterblock}

\begin{posterblock}{Temel Özellikler}
\centering\small
\featcard{myoraOrange}{Y}{Yemek Yakalama}{GPT-4o Vision ile fotoğraftan kalori ve makro besin analizi.}
\hfill
\featcard{myoraPurple}{S}{Sesli Günlük}{NLP ile öğün, aktivite ve ruh hali sesli kaydı.}\\[0.25cm]
\featcard{myoraGreen}{B}{Belge YZ}{Kan tahlili, reçete ve laboratuvar raporlarının OCR + YZ analizi.}
\hfill
\featcard{myoraBlue}{T}{Takviye Motoru}{RDA açıkları ve biyobelirteçlere göre kişisel plan.}\\[0.25cm]
\featcard{myoraCyan}{A}{Analitik}{Haftalık trendler, makro dağılımı ve oruç geçmişi.}
\hfill
\featcard{myoraRed}{F}{Aile Takibi}{Aile bireyleri için ilaç ve uyum izlemi.}
\end{posterblock}

\begin{posterblock}{Metodoloji}
\small
\begin{enumerate}\setlength{\itemsep}{0.2em}
  \item \textbf{Veri Toplama Katmanı}\,--\,Capacitor eklentileri yerel kamera, mikrofon ve dosya seçici API'lerini tek bir çok platformlu arayüzde soyutlar.
  \item \textbf{YZ Analiz Hattı}\,--\,Yakalanan her eser (görüntü / ses transkripsiyonu / PDF) yapılandırılmış sistem istemleriyle \texttt{OpenAI GPT-4o} API'sini çağıran bir Supabase \emph{Edge Function}'a (Deno) yönlendirilir ve tiplenmiş JSON döner.
  \item \textbf{Kişiselleştirme Motoru}\,--\,TypeScript tabanlı RDA hesaplayıcı (NIH/WHO tabloları) yaş, cinsiyet, kilo, aktivite ve VKİ'ye göre mikro besin hedeflerini ayarlar; kan tahlili sonuçları bir biyobelirteç eşleme tablosu üzerinden ilgili RDA değerini çarpar.
  \item \textbf{Durum Yönetimi}\,--\,TanStack React Query, API verilerini kademeli bayatlama süreleriyle önbelleğe alır (10\,sn gerçek zamanlı / 30\,sn dinamik / 5\,dk statik); iyimser güncellemeler yazma Supabase'e ulaşana kadar arayüzü duyarlı tutar.
\end{enumerate}
\end{posterblock}

\begin{posterblock}{Sistem Mimarisi}
\centering
\begin{tikzpicture}[
  node distance=0.55cm,
  every node/.style={font=\small},
  arr/.style={-{Latex[length=3mm]}, line width=1pt, draw=myoraMuted}
]
  \tikzset{layerbase/.style={rounded corners=8pt, line width=1pt,
    minimum width=10cm, minimum height=1.55cm, align=center, inner sep=5pt}}
  \node[layerbase, draw=myoraBlue!60, fill=myoraBlue!10]
    (ui)  {\textbf{İstemci (Android / iOS / Web)}\\React 18 + TS \textbullet\ Tailwind \textbullet\ shadcn/ui \textbullet\ Recharts};
  \node[layerbase, draw=myoraCyan!60, fill=myoraCyan!10, below=of ui]
    (cap) {\textbf{Capacitor 7 Köprüsü}\\Kamera \textbullet\ Mikrofon \textbullet\ Dosya \textbullet\ Haptics \textbullet\ Bildirim};
  \node[layerbase, draw=myoraGreen!60, fill=myoraGreen!10, below=of cap]
    (rq)  {\textbf{React Query Önbelleği}\\iyimser güncelleme \textbullet\ staleTime katmanları};
  \node[layerbase, draw=myoraPurple!60, fill=myoraPurple!10, below=of rq]
    (sb)  {\textbf{Supabase Bulutu}\\Auth \textbullet\ PostgreSQL (RLS) \textbullet\ Storage \textbullet\ Edge Functions};
  \node[layerbase, draw=myoraOrange!60, fill=myoraOrange!10, below=of sb]
    (ai)  {\textbf{Edge Functions (Deno)}\\\texttt{analyze-health-data} \textbullet\ \texttt{generate-chat-response} \textbullet\ \texttt{generate-health-profile}};
  \node[layerbase, draw=myoraRed!60, fill=myoraRed!10, below=of ai]
    (llm) {\textbf{OpenAI GPT-4o \& RAG}\\\texttt{gpt-4o} \textbullet\ \texttt{text-embedding-3-small} \textbullet\ pgvector};

  \draw[arr] (ui)  -- (cap);
  \draw[arr] (cap) -- (rq);
  \draw[arr] (rq)  -- (sb);
  \draw[arr] (sb)  -- (ai);
  \draw[arr] (ai)  -- (llm);
  \draw[arr, dashed, bend left=28, myoraPurple]
    (llm.east) to node[right,text=myoraMuted,align=left,font=\scriptsize]
      {tiplenmiş\\JSON} (ui.east);
\end{tikzpicture}

\vspace{0.1cm}
\small Tüm tablolarda Satır Düzeyi Güvenlik (RLS) politikaları \texttt{user\_id}\,$\rightarrow$\,\texttt{profiles.id} zorunluluğuyla uygulanır.
\end{posterblock}

\end{column}

% ============================================================
% COLUMN 2
% ============================================================
\begin{column}{0.34\textwidth}

\begin{posterblock}{Yapay Zekâ Analiz Hattı}
\centering
\begin{tikzpicture}[
  >={Latex[length=3mm]},
  pbase/.style={
    rounded corners=8pt, line width=1pt,
    minimum width=3.4cm, minimum height=1.9cm, text width=3.2cm,
    align=center, inner sep=4pt,
    font=\footnotesize\bfseries
  },
  arr/.style={->, line width=1.1pt, draw=myoraMuted}
]
  % Row 1: 3 boxes left-to-right
  \node[pbase, draw=myoraBlue!60,   fill=myoraBlue!12,   text=myoraBlue]                         (s1) {Yakalama\\{\tiny\color{myoraMuted}\mdseries resim / ses / PDF}};
  \node[pbase, draw=myoraCyan!60,   fill=myoraCyan!12,   text=myoraCyan,   right=0.45cm of s1]   (s2) {base64\\{\tiny\color{myoraMuted}\mdseries istemci kodlama}};
  \node[pbase, draw=myoraGreen!60,  fill=myoraGreen!12,  text=myoraGreen,  right=0.45cm of s2]   (s3) {Edge Fn\\{\tiny\color{myoraMuted}\mdseries Deno}};
  % Row 2: 3 boxes right-to-left (below row 1)
  \node[pbase, draw=myoraPurple!60, fill=myoraPurple!12, text=myoraPurple, below=1.0cm of s3]    (s4) {GPT-4o\\{\tiny\color{myoraMuted}\mdseries çok modlu}};
  \node[pbase, draw=myoraOrange!60, fill=myoraOrange!12, text=myoraOrange, below=1.0cm of s2]    (s5) {Ayrıştırma\\{\tiny\color{myoraMuted}\mdseries tiplenmiş JSON}};
  \node[pbase, draw=myoraRed!60,    fill=myoraRed!12,    text=myoraRed,    below=1.0cm of s1]    (s6) {Kaydet\\{\tiny\color{myoraMuted}\mdseries Supabase + RQ}};

  \draw[arr] (s1) -- (s2);
  \draw[arr] (s2) -- (s3);
  \draw[arr] (s3) -- (s4);
  \draw[arr] (s4) -- (s5);
  \draw[arr] (s5) -- (s6);
\end{tikzpicture}

\vspace{0.3cm}
\small
Hat, istemci taraflı bir \textbf{yedek mekanizması} içerir: Edge Function 45 sn'yi aşarsa kural-tabanlı tahmini değerler yerel olarak üretilip kaydedilir, böylece kullanıcı deneyimi ağ gecikmesine takılı kalmaz.
\end{posterblock}

\begin{posterblock}{Sağlık Skoru Modeli}
\small
Skor bileşeni, $S \in [0,100]$ aralığında gerçek zamanlı bileşik bir skor üretir:
\begin{equation*}
  \boxed{\; S \;=\; 0{,}30\,N \;+\; 0{,}15\,H \;+\; 0{,}30\,A \;+\; 0{,}25\,C \;}
\end{equation*}
Her bileşen günün verisinden türetilir:

\vspace{0.15cm}
\begin{tikzpicture}[every node/.style={font=\small, text=myoraInk}]
  \node[rounded corners=8pt, draw=myoraOrange!55, fill=myoraOrange!10, line width=1pt,
        minimum width=8.0cm, minimum height=1.55cm, align=left, inner sep=7pt, text width=7.6cm]
    (n1) {\textcolor{myoraOrange}{\bfseries N\ Beslenme (\%30)}\\
          {\scriptsize kalori\_dengesi\,$\times$\,0{,}6 + protein\_\%\,$\times$\,0{,}4}};
  \node[rounded corners=8pt, draw=myoraBlue!55, fill=myoraBlue!10, line width=1pt,
        minimum width=8.0cm, minimum height=1.55cm, align=left, inner sep=7pt, text width=7.6cm,
        right=0.35cm of n1]
    (n2) {\textcolor{myoraBlue}{\bfseries H\ Hidrasyon (\%15)}\\
          {\scriptsize su\_ml\,$\div$\,su\_hedefi\,$\times$\,100}};
  \node[rounded corners=8pt, draw=myoraGreen!55, fill=myoraGreen!10, line width=1pt,
        minimum width=8.0cm, minimum height=1.55cm, align=left, inner sep=7pt, text width=7.6cm,
        below=0.3cm of n1]
    (n3) {\textcolor{myoraGreen}{\bfseries A\ Aktivite (\%30)}\\
          {\scriptsize $\min(\text{aktivite\_sayısı}\times 33,\ 100)$}};
  \node[rounded corners=8pt, draw=myoraRed!55, fill=myoraRed!10, line width=1pt,
        minimum width=8.0cm, minimum height=1.55cm, align=left, inner sep=7pt, text width=7.6cm,
        below=0.3cm of n2]
    (n4) {\textcolor{myoraRed}{\bfseries C\ Tutarlılık (\%25)}\\
          {\scriptsize günlük\_kategori\,$\div$\,4\,$\times$\,100}};
\end{tikzpicture}

\vspace{0.2cm}
\small Formül, \texttt{client/src/components/HealthDashboard.tsx} içinde birebir bu ağırlıklarla uygulanır; bileşen skorları \texttt{ScoreBreakdownModal} üzerinden kullanıcıya açıkça sunulur.
\end{posterblock}

\begin{posterblock}{Kişiselleştirilmiş RDA Motoru}
\small
Takviye modülü, NIH/WHO Diyet Referans Alımlarını kullanarak 18 besin maddesindeki açıkları hesaplar. Aşağıda orta aktiviteli 30 yaşında bir erkek için örnek 7 günlük ortalama karşılanma gösterilmiştir; \textcolor{myoraGreen}{\bfseries yeşil kesikli çizgi} $=$ RDA hedefi.

\vspace{0.2cm}
\centering
\rdabar{Çinko}{62}{myoraRed}\\[0.1cm]
\rdabar{Protein}{88}{myoraOrange}\\[0.1cm]
\rdabar{C Vit.}{74}{myoraYellow}\\[0.1cm]
\rdabar{Ca}{58}{myoraRed}\\[0.1cm]
\rdabar{B12 Vit.}{91}{myoraGreen}\\[0.1cm]
\rdabar{Mg}{66}{myoraOrange}\\[0.1cm]
\rdabar{D3 Vit.}{45}{myoraRed}\\[0.1cm]
\rdabar{Omega-3}{52}{myoraRed}\\[0.1cm]
\rdabar{Demir}{71}{myoraYellow}\\[0.1cm]
\rdabar{Lif}{80}{myoraOrange}

\vspace{0.2cm}
\raggedright
Her açık için doz, sıklık, günün en iyi zamanı, yiyecek eşleşmesi ve ilaç etkileşim uyarılarını içeren sıralı bir \textbf{TakviyePlanı} üretilir. Kan tahlili biyobelirteçleri temel RDA üzerinde çarpan görevi görür (ör. düşük ferritin\,$\Rightarrow$\,demir RDA\,$\times\,1{,}5$).
\end{posterblock}

\begin{posterblock}{Temel Veritabanı Varlıkları}
\centering
\begin{tikzpicture}[
  ent/.style={
    rounded corners=6pt, draw=myoraBlue!55, fill=myoraBlue!8,
    line width=0.9pt, minimum width=4.8cm, minimum height=1.3cm,
    align=center, font=\small\ttfamily
  },
  elink/.style={-{Latex[length=2.5mm]}, draw=myoraMuted, line width=0.7pt}
]
  \node[ent] (p) {profiles};
  \node[font=\tiny, below=0.05cm of p, text=myoraMuted] {kimlik \& onboarding};

  \node[ent, below left=2.1cm and 2.4cm of p]   (f) {food\_entries};
  \node[ent, below=2.1cm of p]                  (a) {activity\_entries};
  \node[ent, below right=2.1cm and 2.4cm of p]  (d) {daily\_health\_logs};

  \node[ent, below=1.1cm of f] (v) {voice\_entries};
  \node[ent, below=1.1cm of a] (b) {blood\_tests};
  \node[ent, below=1.1cm of d] (h) {health\_documents};

  \node[ent, below=1.1cm of v] (s) {supplement\_intake};
  \node[ent, below=1.1cm of b] (m) {family\_members};
  \node[ent, below=1.1cm of h] (u) {user\_streaks};

  \foreach \child in {f,a,d}
    \draw[elink] (p) -- (\child);
\end{tikzpicture}

\vspace{0.25cm}
\small Toplam 32 tablo; hepsinde veritabanı katmanında uygulanan \textbf{Satır Düzeyi Güvenlik} (RLS) politikalarıyla \texttt{user\_id\,($\rightarrow$\,profiles.id)} alanı bulunur. Ek olarak RAG için \texttt{scientific\_sources} ve \texttt{scientific\_source\_chunks} (pgvector embedding) tabloları mevcuttur.
\end{posterblock}

\end{column}

% ============================================================
% COLUMN 3
% ============================================================
\begin{column}{0.32\textwidth}

\begin{posterblock}{Sonuçlar ve Gelecek Çalışmalar}
\small
\textbf{Gerçekleştirilenler:}
\begin{itemize}\setlength{\itemsep}{0.15em}
  \item[\textcolor{myoraGreen}{\checkmark}] Manuel girişin yerini alan \textbf{uçtan uca yapay zekâ gıda analizi} (GPT-4o vision).
  \item[\textcolor{myoraGreen}{\checkmark}] Kan tahlili entegrasyonlu \textbf{kişiselleştirilmiş takviye motoru}.
  \item[\textcolor{myoraGreen}{\checkmark}] Tek kod tabanından \textbf{çok platform} (Android + iOS hedefi + Web) desteği.
  \item[\textcolor{myoraGreen}{\checkmark}] Capacitor Preferences + React Query önbelleğiyle \textbf{çevrimdışı öncelikli} mimari.
  \item[\textcolor{myoraGreen}{\checkmark}] \texttt{pgvector} üzerinde kaynaklı atıf veren \textbf{RAG sohbet} katmanı.
  \item[\textcolor{myoraGreen}{\checkmark}] Edge Function başarısız olduğunda kural-tabanlı \textbf{yedek yanıt} ile servis sürekliliği.
\end{itemize}

\textbf{Gelecek:}
\begin{itemize}\setlength{\itemsep}{0.15em}
  \item[\textcolor{myoraOrange}{$\rightarrow$}] Apple HealthKit ve Google Health Connect ile \textbf{canlı giyilebilir senkron} (Apple Watch, Galaxy Watch).
  \item[\textcolor{myoraOrange}{$\rightarrow$}] Open Food Facts API üzerinden \textbf{ambalajlı gıda barkod tarayıcısı}.
  \item[\textcolor{myoraOrange}{$\rightarrow$}] Sesli günlükler için \textbf{OpenAI Whisper} transkripsiyon entegrasyonu.
  \item[\textcolor{myoraOrange}{$\rightarrow$}] Kişiselleştirilmiş kalori tahmini için \textbf{boylamsal makine öğrenmesi} modeli.
  \item[\textcolor{myoraOrange}{$\rightarrow$}] Sürekli glukoz monitörü (CGM) entegrasyonu.
  \item[\textcolor{myoraOrange}{$\rightarrow$}] GPT-4o-mini ile maliyet kademeleme ve klinik doğrulama çalışması.
\end{itemize}
\end{posterblock}

\begin{posterblock}{Etkileşim ve Oyunlaştırma}
\small
Myora, günlük tutarlılığı artırmak için \emph{Duolingo tarzı} mekanikler içerir:
\begin{itemize}\setlength{\itemsep}{0.1em}
  \item \textcolor{myoraYellow}{\bfseries Başarımlar}\,--\,seriler, hedefler ve günlük giriş için rozetler.
  \item \textcolor{myoraOrange}{\bfseries Seriler}\,--\,kurtarma bildirimleriyle günlük kontrol serisi.
  \item \textcolor{myoraPurple}{\bfseries XP ve Seviyeler}\,--\,öğün, aktivite ve ses günlükleri XP kazandırır; 5 seviye.
  \item \textcolor{myoraGreen}{\bfseries Günlük Görevler}\,--\,her gün sıfırlanan 5 görev; tam tamamlanma kutlanır.
  \item \textcolor{myoraBlue}{\bfseries Mini Zorluklar}\,--\,hafta indeksine göre dönen 3 haftalık meydan okuma.
\end{itemize}
Tüm mekanikler \texttt{user\_streaks}, \texttt{achievements}, \texttt{user\_achievements} tabloları ve \texttt{challengeAPI} üzerinden kalıcılaştırılır.
\end{posterblock}

\begin{posterblock}{Etik ve Güvenlik}
\small
\begin{itemize}\setlength{\itemsep}{0.1em}
  \item \textbf{RLS} her tabloda \texttt{user\_id} üzerinden aktif; ayrıcalık yükseltme engellenir.
  \item Sistem bir \textbf{tıbbi tanı aracı değil}; klinik kararlar için profesyonel hekimliğe yönlendirir.
  \item Tıbbi belge depoları Supabase Storage'da özel bucket'larda tutulur.
  \item Edge Function'lar OpenAI hatalarında \textbf{bilgi sızdırmaz}, jenerik fallback döner.
  \item Halüsinasyon riskini düşürmek için sohbet, \textbf{pgvector} üzerinden bilimsel kaynaklarla kısıtlı RAG bağlamı kullanır.
\end{itemize}
\end{posterblock}

\begin{posterblock}{Teknik Özellikler}
\small
\begin{tabularx}{\linewidth}{@{}lY@{}}
\toprule
\textbf{Katman} & \textbf{Bileşen} \\
\midrule
UI         & React 18 \textbullet\ TypeScript 5.8 \textbullet\ Vite 5 \textbullet\ Tailwind CSS \textbullet\ shadcn/ui \textbullet\ Recharts \\
Durum      & TanStack Query 5 \textbullet\ React Context \textbullet\ Zod \\
Mobil      & Capacitor 7 \textbullet\ Camera \textbullet\ Mic \textbullet\ Local Notifications \textbullet\ Preferences \textbullet\ Haptics \\
Arka uç    & Supabase Auth \textbullet\ PostgreSQL 15+ \textbullet\ RLS \textbullet\ Storage \textbullet\ Deno Edge Functions \\
Yapay zekâ & OpenAI GPT-4o (vision + text) \textbullet\ text-embedding-3-small \textbullet\ pgvector RAG \\
Kalite     & ESLint 9 \textbullet\ strict TS \textbullet\ lazy-loaded rotalar \textbullet\ ErrorBoundary \\
\bottomrule
\end{tabularx}

\vspace{0.3cm}
\centering
\tagpill{myoraBlue}{React 18}\ \tagpill{myoraCyan}{TypeScript}\ \tagpill{myoraGreen}{Vite 5}\ \tagpill{myoraOrange}{Tailwind}\\[0.12cm]
\tagpill{myoraPurple}{Supabase}\ \tagpill{myoraRed}{PostgreSQL}\ \tagpill{myoraBlue}{Deno}\ \tagpill{myoraGreen}{RLS}\\[0.12cm]
\tagpill{myoraOrange}{GPT-4o}\ \tagpill{myoraPurple}{Vision API}\ \tagpill{myoraCyan}{RAG}\ \tagpill{myoraGreen}{Zod}\\[0.12cm]
\tagpill{myoraRed}{Capacitor}\ \tagpill{myoraBlue}{iOS/Android}\ \tagpill{myoraYellow}{React Query}\ \tagpill{myoraGreen}{pgvector}
\end{posterblock}

\begin{posterblock}{Derleme ve Kaynak}
\small
\begin{itemize}\setlength{\itemsep}{0.1em}
  \item \textbf{Ana dosya:} \texttt{poster/myora\_academic\_poster.tex}
  \item \textbf{Derleyici:} pdfLaTeX veya Overleaf
  \item \textbf{Gerekli paketler:} \texttt{beamerposter}, \texttt{tikz}, \texttt{babel} (Turkish), \texttt{xcolor}, \texttt{tabularx}, \texttt{amsmath}
\end{itemize}

\vspace{0.2cm}
\textbf{Repo:} \href{https://github.com/seyityorgun499-cyber/wellness-vision-coach}{github.com/seyityorgun499-cyber/wellness-vision-coach}\\[0.15cm]
\textbf{Veri kaynakları:} NIH / WHO Dietary Reference Intakes tabloları, RAG için PubMed ve akademik derleme bağlantıları (\texttt{scripts/rag/seed-sources.ts} içinde).
\end{posterblock}

\end{column}

\end{columns}

\vfill
\begin{center}
\small\color{myoraMuted}
© 2026 Myora Health\ \ \textbullet\ \ React \& Capacitor istemci\ \ \textbullet\ \ Supabase + Deno Edge Functions\ \ \textbullet\ \ OpenAI GPT-4o\ \ \textbullet\ \ \texttt{github.com/seyityorgun499-cyber/wellness-vision-coach}
\end{center}

\end{frame}
\end{document}
